\section{Analysis design and statistical methods}
\label{sec:analysis_design}

This section defines the operating points, estimators, and uncertainty
procedures used for the experiment sets in Table~\ref{tab:simulation_sets}.

\subsection{Evaluation cohorts, scoring procedure, and operating points}
\label{sec:scoring_contract}

The primary selected-system endpoint uses the 43-system evaluation cohort
defined in Section~\ref{sec:data_partitions}. None was used for detector
fitting or validation, detector-model selection, threshold calibration, or
exploratory feature-analysis development. The 20 systems assigned to that
exploratory role are excluded from this denominator.

Two disjoint catalogue-negative samples are drawn as whole spatial groups
from the detector's negative test population. The threshold sample contains
20,002 rows in 1,768 spatial groups. A second sample contains 12,512 rows in
1,089 spatial groups and is used only to measure the achieved FPR without
modifying the calibrated threshold.

For the primary analyses, the selected-system endpoint, \(G_{\rm rec}\),
\(P_{\rm noise}\), and \(P_{\rm src}\) use the same epoch-160
EfficientNetV2-S checkpoint, preprocessing, and scoring procedure. The
network is in evaluation mode:
all batch-normalisation layers use stored running statistics and dropout is
inactive. Inference uses float32 operations.
Logits define the operating point because extreme float32 sigmoid values can
saturate numerically. The displayed
\(p=\operatorname{sigmoid}(\ell)\) is not treated as a calibrated
probability. Execution-environment and batching details are reported in
Appendix~\ref{app:technical_validation}.

For \(N_0\) calibration negatives with native float32 logits sorted in
descending order
\(s_{(1)}\geq\cdots\geq s_{(N_0)}\), a target FPR \(\alpha\) permits
\(k=\lfloor\alpha N_0\rfloor\) exceedances. We choose
\begin{equation}
\begin{aligned}
 t_\alpha
 &=\operatorname{nextafter}_{32}\!\left(s_{(k+1)},+\infty\right),\\
 D_i(t_\alpha)&=\mathbb{I}(\ell_i\geq t_\alpha).
\end{aligned}
\label{eq:common_threshold}
\end{equation}
With \(N_0=20\,002\), \(\alpha=10^{-3}\), and \(k=20\),
this gives \(t_{10^{-3}}\simeq1.572\) and 20 calibration exceedances. The
full-precision float32 value is reported in
Appendix~\ref{app:technical_validation}. The threshold is applied unchanged
to every primary detector-response analysis; the feature probes do not use it.

\subsection{Detector-checkpoint sensitivity}
\label{sec:checkpoint_sensitivity_methods}

The primary detector is the seed-42 epoch-160 checkpoint trained with reviewed
and additional catalogue candidates. The sensitivity analysis selects the
highest-validation-AP checkpoint from each of the 12 training runs, retaining the earlier epoch
under ties. Each checkpoint receives its own threshold from the same
calibration sample using Equation~\ref{eq:common_threshold}, and that
threshold is applied to the common evaluation data.
The target \(10^{-3}\) defines the empirical calibration rate.
Checkpoint-specific estimates and intervals condition on \(t_m\).

The pixels, estimators, and dependence-aware resampling procedures are
unchanged across checkpoints; only model weights, scores, and thresholds
vary. The operational \(G_{\rm domain}\) and \(G_{\rm match}\) probes
characterize only the designated primary checkpoint and are not repeated
across the panel. The panel tests how thresholded detector responses change
across trained models; the feature-space analysis is performed only for the
primary detector. The 12 checkpoints are correlated variants within one
detector design \citep{Schaefer2018,Canameras2024}. We report both
validation-selected and terminal-epoch checkpoints; comparisons between
these checkpoint policies are exploratory.

\subsection{Detector-performance and recovery estimators}
\label{sec:performance_estimators}

For a labelled population at a stated threshold, recall is
\({\rm TP}/({\rm TP}+{\rm FN})\), FPR is
\({\rm FP}/({\rm FP}+{\rm TN})\), and precision is
\({\rm TP}/({\rm TP}+{\rm FP})\). ROC AUC measures ranking across all
thresholds, with half credit for ties. The reported precision--recall summary
is average precision. The Brier score is
\begin{equation}
 B = \frac{1}{n}\sum_{i=1}^{n}(p_i-y_i)^2 .
\label{eq:brier_score}
\end{equation}
Validation-positive metrics use noisy reviewed and additional catalogue
candidate labels and are development diagnostics, not estimates on confirmed
systems.

For a named injection population \(Q\), recovery at threshold \(t\) is
\begin{equation}
 \widehat C_Q(t)
 = \frac{1}{|Q|}\sum_{i\in Q}\mathbb{I}(\ell_i\geq t).
\label{eq:conditional_recovery}
\end{equation}
This quantity is conditional on the target population, source and mass-model
priors, observing-condition grid, visibility policy, detector, and threshold
specified for \(Q\). We call this quantity conditional recovery rather than
survey-population completeness.
Because \(G_{\rm rec}\) imposes no independent detectability or visibility
condition, its recovery fraction combines configurations that may not be
physically detectable with the finder's response; it is not an estimate of
finder response conditional on detectability alone.

For \(x\) row-level successes among \(n\) rows, the 95 per cent Wilson score
interval \citep{Wilson1927} is
\begin{equation}
\begin{aligned}
 p_{\rm W} &=
 \frac{\widehat p+z^2/(2n)}{1+z^2/n},\\
 \Delta_{\rm W} &=
 \frac{z}{1+z^2/n}
 \sqrt{\frac{\widehat p(1-\widehat p)}{n}+\frac{z^2}{4n^2}},
 \\
 {\rm CI}_{\rm W} &= [p_{\rm W}-\Delta_{\rm W},p_{\rm W}+\Delta_{\rm W}],\\
 \widehat p &= x/n,\qquad z=\Phi^{-1}(0.975)\simeq1.96 .
\end{aligned}
\label{eq:wilson_interval}
\end{equation}
We apply this row-level interval to the achieved-FPR endpoint. Because nearby
catalogue rows can be dependent, that endpoint also receives a 10,000-draw
spatial-group interval: each draw assigns one independent mean-one
exponential multiplier to each spatial group and takes the weighted ratio of
exceedances to rows. The Wilson and multiplier intervals correspond to
different dependence summaries. The 43 selected confirmed systems are
reported as an exact finite cohort without a binomial or population interval.

For \(G_{\rm rec}\), let \(D_{chr}\) be the threshold indicator for target
\(h\) in retained observing cell \(c\) at radius \(r\). With nine cells, 500
targets in each cell, and 11 radii per target, the primary point estimator is
\begin{equation}
\widehat C_{\rm rec}
=\frac{1}{9}\sum_{c=1}^{9}\frac{1}{500}\sum_{h=1}^{500}
  \frac{1}{11}\sum_{r=1}^{11}D_{chr}.
\label{eq:grec_balanced_recovery}
\end{equation}
Thus cells, targets within cells, and the 11 radius nodes receive equal
weight. This equals the raw row mean because the design is balanced; it is
not a survey-population weighting.

Sampling uncertainty uses 10,000 exponential-multiplier draws at the
injection-target spatial-group level. A group receives one global multiplier
per draw, and that same multiplier follows all of its targets wherever they
occur across cells and follows every radius of each target. Within each cell,
the weighted target mean is formed as a ratio. To remove finite-draw ratio
bias without changing the joint covariance, the draw distribution for every
cell--radius entry is recentered to its exact equal-target point estimate.
The overall 95 per cent interval is the 2.5th and 97.5th percentiles of the
recentered, equal-cell and equal-radius draws.

The radius curve averages cells and targets at each radius. Its unstudentised
simultaneous band uses
\begin{equation}
\begin{aligned}
 q_{0.95}
 &=\operatorname{quantile}_{0.95}
 \left(
 \left\{\max_r\left|\widehat C^{*(b)}_r-\widehat C_r\right|\right\}
 _{b=1}^{10000}\right),\\
 {\cal B}_r&=[\widehat C_r-q_{0.95},\,\widehat C_r+q_{0.95}] .
\end{aligned}
\label{eq:grec_simultaneous_band}
\end{equation}
The band is not clipped to \([0,1]\). Every cell--radius estimate has 500
distinct targets; the resampling unit, rather than the injection count,
defines independent support.

\subsection{Linear probes and grouped evaluation}
\label{sec:linear_probe_methods}

A linear probe consists of fold-local standardisation followed by
\(L_2\)-regularised logistic regression on features computed with encoder
weights held fixed. Every probe uses \(C=1\), the \texttt{lbfgs} solver,
\texttt{max\_iter}=5000, estimator seed 42, and scikit-learn 1.9.1. The
weighted mean and variance used by \texttt{StandardScaler} are fitted only on
the training folds, and the same training weights are supplied to the
logistic-regression fit.

Dependency groups are defined without examining feature values. They are the
transitive union of shared physical group, NSIDE\(=128\) NESTED HEALPix cell
\citep{Gorski2005}, within-population spatial group, reference identity,
injection-target identity and underlying-scene identity, source, lens, PSF, renderer and
render-parameter identities, random-draw identity, and base or final pixel
hash. Every resulting dependency component remains wholly in one of five
folds; score-blind fold allocation uses seed 42. This closes indirect links,
including a reference geographically linked to an injection target in a
different nominal pair.

Across the full set of rows eligible to fit an axis, the two fitting classes
receive equal total base weight. Within a class, every dependency component receives equal total
weight, divided equally over its rows. Each fold fit uses the corresponding
weights of its retained training components without fold-specific
renormalisation. For fitting-roster size \(N\), class \(y\) represented by
\(K_y\) components, and \(n_{c,y}\) class-\(y\) rows in component \(c\), the
weight is
\begin{equation}
 w_i=\frac{N}{2K_y n_{c,y}}.
\label{eq:probe_weights}
\end{equation}
For \(G_{\rm domain}\), the weights sum to 6,379, with 3,189.5 assigned to
each class. Consequently, repeated radii from one connected target do
not receive eleven times the training weight of a singleton component. The
same full-roster component-equal, class-balanced weights define the primary
pooled out-of-fold (OOF) ROC AUC. A row-weighted OOF AUC is retained only as
a diagnostic.
ROC AUC uses half credit for tied positive--negative scores.

The domain probe compares all 43 selected confirmed systems with the exact
6,336-row \(G_{\rm domain}\) subset in 576 injection targets. The common
dependency graph contains 561 components: 43 contain observed rows, 520
contain simulated rows, and two contain both. There are zero component
crossings over the five folds. The comparison uses four feature sets. The
operational representation \(z(x)\) is the 1,280-component output after the
primary detector's global-average-pooling layer and before its dropout and
linear classifier. The existing lens head evaluates
\begin{equation}
 h(x)=w_{\rm lens}^{\mathsf T}z(x)+b .
\label{eq:fixed_head}
\end{equation}
The operational domain probe instead fits fold-specific regularized linear
readouts of fold-standardized \(z(x)\) to observed-versus-injected labels,
using training rows only. This examines the representation that feeds the
finder without restricting the diagnostic to its learned lens-classification
direction. The existing lens head is reported as a fixed-score baseline
alongside the four fitted feature protocols. Both use the
detector-native full \(101\times101\) \(g,r,z\) input with the trained mapping
defined in Section~\ref{sec:preprocessing}.

The Stein representation is the 2,048-component global-average-pooled output
of its self-supervised ResNet-50 before the projection head
\citep{Stein2022}. Its encoder was trained with a contrastive objective on
unlabelled Legacy Survey images. Object-level overlap between those
pretraining images and the present cutouts is unknown, so the Stein result is
interpreted as feature-protocol sensitivity. The generic ImageNet
representation is the 1,280-component global-average-pooled output of a
separate EfficientNetV2-S, with different weights from the lens-label-conditioned
operational detector. The fourth representation contains 18 fixed pixel
statistics. Each representation includes its own preprocessing. The two
external deep encoders use the Legacy Survey display mapping adopted by
\citet{Stein2022}. For nanomaggy
channels \(x_c\), let \(u_c=a_cx_c+m\), with
\((a_g,a_r,a_z)=(6.0,3.4,2.2)\), \(m=0.03\), and \(q=20\). The exact mapping is
\begin{equation}
\begin{aligned}
 I&={1\over3}\sum_{c\in\{g,r,z\}}\max(0,u_c),\\
 F(I)&={\operatorname{asinh}(qI)\over\sqrt q\,I_{\rm safe}},\\
 (R,G,B)&=\operatorname{clip}_{[0,1]}\!\left[(u_z,u_r,u_g)F(I)\right],
\end{aligned}
\label{eq:rgb_preprocessing}
\end{equation}
where \(I_{\rm safe}=I+10^{-6}\mathbf{1}[I=0]\). The Stein input is
the central \(96\times96\) pixels. The ImageNet input uses all
\(101\times101\) pixels, is bilinearly resized with antialiasing to
\(384\times384\), and receives the standard ImageNet channel normalisation.
The pixel representation records, separately in each \(g,r,z\) band, the
mean, standard deviation, median, 99th percentile, minimum, and maximum.
Probe comparisons therefore diagnose the frozen representation together with
its preprocessing, not encoder weights alone.

\(G_{\rm domain}\) reports dependency-component-equal pooled OOF AUCs as
finite-cohort, model-conditional point estimates. No resampling interval or
significance test is attached. We compare all four representations on the
same rows, folds, component weights, and observed--injected labels. The five
fold-level AUCs show how performance varies across held-out subsets containing
only eight or nine observed systems.

\subsection{Unpaired separation and matched controls}
\label{sec:decomposition_methods}

We distinguish three questions. An unpaired real--injection probe asks whether
the selected populations are separable. A real--bare-target probe measures
separation present before an arc is added. A bare--injected comparison holds
the target fixed and measures how the added arc moves that target in feature
space. We interpret these estimands separately because they answer different
population and paired-movement questions.

The matched-baseline set \(G_{\rm match}\) uses the exact 43 selected confirmed
systems as observed references. They are matched without replacement to the
4,500 \(G_{\rm rec}\) injection targets on observed \(z\)-band magnitude,
catalogue S\'{e}rsic index, axis ratio, \(r\)-band PSF FWHM, and \(r\)-band
depth. No detector score or representation feature enters matching.

For each covariate \(k\), the complete 4,500-target sample supplies the median
\(m_k\) and interquartile range \(s_k\), using all finite values of that
component. We define \(z_{ik}=(x_{ik}-m_k)/s_k\). All five values must be
finite for a target to be eligible; all 4,500 targets meet this requirement,
and no imputation is used. A reference--target edge is allowed only when
\begin{equation}
\begin{aligned}
 |z_{{\rm ref},k}-z_{{\rm target},k}|&\leq2
 &&\text{for every }k,\\
 \left[\sum_k(z_{{\rm ref},k}-z_{{\rm target},k})^2\right]^{1/2}
 &\leq3 .
\end{aligned}
\label{eq:gmatch_calipers}
\end{equation}
Among allowed edges, maximum-cardinality then minimum-total-distance
assignment is used, with SHA-256-sorted identifiers breaking ties. All 43
references are matched to unique targets.

The supplementary material includes the complete full-field reference--target
atlas, the matching and brightness-parent tables, covariate balance summaries,
and the code used to generate the displays.

Each target contributes its exact stored bare cutout and its exact
\(\theta_{\rm E}=1.5\)-arcsec \(G_{\rm rec}\) deterministic expectation image. There is no
rerendering, post-render selection, flux rescaling, imputation, replacement,
or target-brightness floor. The resulting 43 observed--bare--injected
triplets form 43 closed dependency components.

In each representation and fold, the logistic axis is fitted only to the
training-fold observed and bare rows, with observed as the positive direction;
injected rows fit neither the axis nor its scale. Let \(a_{jf}^{x}\) denote a
held-out decision-function score for target \(j\), fold \(f\), and arm
\(x\in\{{\rm bare},{\rm inj}\}\). We standardise both arms using the weighted
mean \(\mu_f^{\rm bare}\) and population standard deviation
\(\sigma_f^{\rm bare}\) of the training-fold bare scores:
\begin{equation}
\begin{aligned}
\widetilde a_{jf}^{\,x}
&=\frac{a_{jf}^{x}-\mu_f^{\rm bare}}{\sigma_f^{\rm bare}},\\
M_j&=\widetilde a_{jf}^{\,{\rm inj}}-\widetilde a_{jf}^{\,{\rm bare}},
\qquad
\widehat\Delta_{\rm match}=\frac{1}{43}\sum_{j=1}^{43}M_j .
\end{aligned}
\label{eq:gmatch_movement}
\end{equation}
Here \(M_j\) is an equal-target movement measured in training-bare
standard-deviation units toward the observed-reference side of a cross-fitted
axis.

Pooled observed--bare, observed--injected, and injected--bare AUCs, with half
credit for ties, summarize the same fitted axis. These point summaries
condition on the frozen score-blind assignment; no formal interval or
significance test is attached. The matching set contains five measured
central-object and observing quantities; complete photometric-redshift and
source-magnitude covariates are unavailable, and the observed central-object
quantities can include lensed-source light.

\subsection{Fixed-target brightness and source-photon response}
\label{sec:bounded_control_methods}

\(C_{\rm bright}\) uses equation~\eqref{eq:brightness_control} on 43 fixed
targets under the detector's trained 20--32-pixel annulus and clip-10
preprocessing. For each target and magnitude, photon indicators are first
averaged over the eight keyed draws; the resulting target-level values are
then averaged equally over targets. Primary summaries are recovery and
native-logit response over magnitude; photon-minus-expectation is the paired
secondary contrast. If \(\widehat R_d\) is the equal-target photon recovery at
one magnitude for draw index \(d\), the finite-draw diagnostic is
\begin{equation}
 {\rm MCSE}_{\rm bright}
 =\frac{{\rm sd}(\widehat R_1,\ldots,\widehat R_8)}{\sqrt{8}} .
\label{eq:brightness_mcse}
\end{equation}
No population interval or visibility admission rule is attached. Source-only
in-stamp flux and peak intensity are diagnostics rather than calibrated arc
S/N.

\subsection{Arc-level Poisson experiment}
\label{sec:poisson_methods}

\(P_{\rm noise}\) compares all 49,500 stored \(G_{\rm rec}\) expectations
with eight modelled source-photon realizations per scene at
\(g=150\ {\rm e^{-}\,nmgy^{-1}}\), using the common threshold. The target
and expected arc remain fixed; only source-photon realization changes
(Section~\ref{sec:noise_model}).

Let \(D^{\rm E}_{chr}\) denote the deterministic expectation decision and
\(D^{\rm M}_{chrd}\) the modelled-photon decision for draw \(d\). The primary
modelled-photon-minus-expectation estimand is
\begin{equation}
\widehat\Delta_{\rm noise}
=\frac{1}{9}\sum_{c=1}^{9}\frac{1}{500}\sum_{h=1}^{500}
 \frac{1}{11}\sum_{r=1}^{11}
 \left\{\frac{1}{8}\sum_{d=1}^{8}D^{\rm M}_{chrd}-D^{\rm E}_{chr}\right\}.
\label{eq:pnoise_risk_difference}
\end{equation}
Photon draws are therefore averaged within deterministic expectation scene
before radii are averaged within target; targets and retained cells then
receive equal weight. Repeated photon draws never increase a target's
scientific weight.

The sampling interval uses 10,000 target-spatial-group exponential-multiplier
draws with the same global grouping, equal-cell weighting, and cell--radius
recentering as \(G_{\rm rec}\). Its uncertainty is conditional on the
realized target and source construction and does not include gain-model or
threshold-calibration uncertainty. The radius-specific effects receive the
same maximum-absolute-deviation simultaneous-band construction as
equation~\eqref{eq:grec_simultaneous_band}. Gained and lost threshold
crossings are reported as descriptive probabilities over the complete
scene--draw rows.

For scene \(i\), let \(s_i^2\) be the sample variance of its eight modelled
threshold indicators and let \(w_i=1/49\,500\), which is the final
equal-cell, equal-target, equal-radius weight under the balanced design. The
finite-draw contribution is
\begin{equation}
 {\rm MCSE}_{\rm noise}
 =\left(\sum_{i=1}^{49\,500}w_i^2\frac{s_i^2}{8}\right)^{1/2}.
\label{eq:pnoise_mcse}
\end{equation}
The estimate applies to the specified stochastic source-photon model, whose
\(g=150\) conversion and post-coadd construction define the intervention.

\subsection{Paired source estimators}
\label{sec:paired_source_estimators}

For pair \(j\) and photon draw \(d\), let \(D^{\rm ser}_{jd}\) and
\(D^{\rm str}_{jd}\) indicate whether the corresponding photon-realized images
from the single-S\'{e}rsic and structured constructions exceed the common
logit threshold. The underlying source constructions are deterministic, but
the scored rows are draw-specific images. The primary estimand is the
S\'{e}rsic-minus-structured photon-averaged risk difference,
\begin{equation}
\widehat\Delta_{\rm src}
=\frac{1}{1584}\sum_{j=1}^{1584}\frac{1}{8}\sum_{d=1}^{8}
 \left(D^{\rm ser}_{jd}-D^{\rm str}_{jd}\right).
\label{eq:psrc_risk_difference}
\end{equation}
Operationally, draw differences are first averaged within pair; pairs are
then averaged within each of the nine fixed target cells and cells are
averaged equally. The expression above is identical because every cell has
the same number of pairs, every cell--radius block contains all 16 source
strata, every selected source is used 11 times, and all base analysis weights
are one.

The 95 per cent interval is a 10,000-draw multiway exponential-multiplier
interval. Each draw assigns one global mean-one exponential multiplier to
each of the 1,219 target spatial groups and applies it to that group's rows in
every cell and radius. It independently assigns one multiplier to each of the
144 empirical sources; source multipliers are normalised to mean one within
each of the 16 fixed source strata. The product of the spatial and source
multipliers weights each pair, after which equal-cell ratio estimates are
averaged. Draws are recentered to \(\widehat\Delta_{\rm src}\), and the 2.5th
and 97.5th percentiles form the reported interval. This procedure preserves
dependence from repeated sources and geographically linked targets. The
resulting interval is a design-based approximation, not an exact
binary-outcome interval.

Finite-photon-draw Monte Carlo uncertainty is reported separately. Let
\(\widehat\Delta_d\) be the same equal-cell risk difference computed from
draw index \(d\) alone. Then
\begin{equation}
 {\rm MCSE}_{\rm src}
 =\frac{{\rm sd}(\widehat\Delta_1,\ldots,\widehat\Delta_8)}{\sqrt{8}} .
\label{eq:psrc_mcse}
\end{equation}
The point estimate and multiplier interval condition on the
calibrated threshold. The empirical construction retains measurement
uncertainty, making the estimand a comparison of the complete
source-construction pathways.

\subsection{Interpretation of estimates and intervals}
\label{sec:reporting_methods}

The achieved-FPR, \(G_{\rm rec}\), \(P_{\rm noise}\), and \(P_{\rm src}\)
intervals use the 95 per cent procedures defined above. The selected
43-system endpoint and both feature-probe analyses are finite-cohort point
summaries without population intervals. The four feature protocols form one
finite-cohort comparison. The operational probes characterize the primary epoch-160
checkpoint; the 12-checkpoint panel evaluates thresholded detector-response
sensitivity and does not repeat the feature probes.

Each thresholded estimate conditions on the threshold calibrated
for its checkpoint.
Multiplier intervals describe the specified target/source dependence but do
not incorporate detector-training, threshold-calibration, gain-model, source
reconstruction, or population-transport uncertainty. Finite-photon MCSE is
reported separately from sampling uncertainty. Formal equivalence assessment
would require a predefined margin beyond these estimation summaries.
